\documentclass[slidetop]{beamer}

\usepackage[utf8x]{inputenc}
\usepackage[francais,french]{babel}
\usepackage{fontenc}
\usepackage{graphicx}
\usepackage{url}

\usetheme{Copenhagen}
\usepackage{amssymb}

\newcommand{\coquille}{\textsc{Coquille}}
\newcommand{\coqcode}[1]{\texttt{#1}}

\date{10 novembre 2009}
\title{Rapport du projet \coquille{}}
\author{M1 Informatique Fondamentale ENS Lyon}

\usepackage{default}

\begin{document}
\section{Introduction}
\begin{frame}
\maketitle
\center
\includegraphics[width=0.25\textwidth]{poussin.png}

\end{frame}

\begin{frame}
\tableofcontents

\end{frame}

\section{Présentation}
\begin{frame}
Coquille : Coq User-Interactive Library Learning Expert

Objectifs :
\begin{itemize}
  \item Objectif: automatiser la résolution des exercices de mathématiques des classes préparatoires.
    %% à propos, il faut changer les objectifs
  \item Créer les bibliothéques correspondant aux mathématiques de classes préparatoires dans Coq. 
  \item Rendre Coq accessible aux mathématiciens. % non forcément experts en calcul des constructions et en théorie des types. 
\end{itemize}


\end{frame}



\section{WP Preuves et tactiques}
\begin{frame}
\frametitle{WP Preuves et tactiques} 

Le but du WP preuves:
\begin{itemize}
  \item établir une liste de théorèmes vus en classes préparatoires en: 
  \begin{itemize}
    \item algèbre,
    \item analyse,
    \item topologie,
    \item arithmétique.
  \end{itemize}
  \item formaliser les thèorèmes en Coq pour les prouver.
  \item proposer des tactiques évoluées d'aide aux prouveurs. 
  %des tactiques d'aide ?
\end{itemize}

\end{frame}


\begin{frame}
 
Quelques jolis résultats:
\begin{itemize}
  \item R est indénombrable.
  \item L'axiomatique des réels implique le principe de Markov.
  \item La limite de zeta (2) est pi²/6.
  %zeta(2) est une constante !
  \item Convergence de suites réelles et construction de limites.
\end{itemize}
\end{frame}

\begin{frame}
De nouvelles bibliothèques:
\begin{itemize}
  \item Arithmétique.
  \item Topologie (définition de nombreuses typeclass de base, recherche de définitions les plus générales possibles).
    %ça a changé, ce n'est pas que de la topologie
  \item suites et séries réelles et complexes.
  \item Définition du corps des complexes, résultats de base.
    %on ne parle pas de concepts.
  \item Formalisation des séries entières et développements sur les rayons de convergence.
\end{itemize}
\end{frame}


\begin{frame}
\begin{itemize}
\item Toute la documentation générée par coqdoc est disponible sur le site du projet. 
\item Le respect des conventions de nommage permet une utilisation facilitée.
\end{itemize}
\end{frame}


\section{WP IDE}
\begin{frame}
\frametitle{WP IDE} 
L'objectif du WP IDE est de construire une interface de programmation en Coq plus orientée vers l'aide à la programmation que l'actuel CoqIde.
%% L'AIDE À LA PROGRAMMATION ?
% non, pas "plus orientée vers l'aide" non plus. juste plus élaborée avec plus de fonctionalités, et moins de bug (mais pas pour l'instant)
\end{frame}

\begin{frame}
Ce que Coquille fait et que CoqIde ne fait pas au niveau du code :
\begin{itemize}
    \item Les numéros de ligne.
    \item Le code folding (replier des lignes de code en une seule, pour une preuve par exemple).
    \item Des raccourcis plus "classiques", et personnalisables.
    \item Possibilité de faire "Redo" après des "Undo".
\end{itemize}
\end{frame}

\begin{frame}

Au niveau du langage :
\begin{itemize}
    \item La gestion de Ltac Debug, avec une interface de parcours des résultatsScreenshot
    \item Plusieurs instances de Coqtop, une par onglet ouvert
    \item Un affichage des résultats au choix : classique ou LaTeX-likeScreenshot
    \item L'action "Next/Previous" d'envoi d'une commande est considérée comme une action comme les autres, donc "Undo/Redo" agit dessus
\end{itemize}
 
\end{frame}

\begin{frame}

Ce que CoqIde fait au niveau du code :

\begin{itemize}
    \item Lister les tactiques existantes dans un menu.
    \item Lister les actions disponibles par click droit sur une hypothèse ou un but.
\end{itemize}



\end{frame}


\begin{frame}

Au niveau du langage :

\begin{itemize}
    \item Le Proof Wizard.
%    \item La gestion des Write State / Restore State
    \item La gestion de l'aide.
    \item La gestions de la compilation.
\end{itemize}


\end{frame}



\begin{frame}
\includegraphics[width=1\textwidth]{unicode.png}
%ça ne compile pas ! Ajoutez les images, bordel --> ``attention, des images ont été oubliées'' est plus agréable à lire.
%montrer des choses que ne peut pas faire coquille (quoi ?)
\end{frame}

\begin{frame}
\includegraphics[width=1\textwidth]{ltacdebug.png}
%(ltac debug)
\end{frame}

\section{WP Apprentissage}
\begin{frame}
\frametitle{WP Apprentissage} 
WP Apprentissage
L'objectif de ce groupe de travail est à partir de nombreux exercices, de créer des tactiques spécifiques à ce type d'exercice.
\end{frame}
\begin{frame}

% schéma nécessaire
\begin{itemize}
    \item On utilise un arbre de décision. 
    \item A chaque noeud se trouve une question portant sur les arbres syntaxiques des hypothèses et du but. 
\begin{itemize}
    \item Pour la première question, on ne s'autorise à interroger que les racines.
    \item Ensuite, on s'autorise à interroger les fils des arbres déjà interrogés.
\end{itemize}
\end{itemize}

\end{frame}

\begin{frame}

Les questions sont:
\begin{itemize}
    \item Quel est le token à la racine de l'arbre x?
    \item Est-ce que les arbres x et y ont le même token en racine?
\end{itemize}


Pour choisir la question à poser, 
on calcule la question qui nous donne le meilleur rapport "information apportée selon la définition de von Newmann" sur "nombre de réponses". 
On cherche donc à minimiser la taille de l'arbre (et non sa profondeur),
car c'est d'elle que dépend le temps d'exécution et l'occupation de la mémoire.
\end{frame}
\begin{frame}

Résultats
\begin{itemize}
    \item Le parcours d'arbres et le raffinement des données sont utilisables.
    \item L'algorithme d'apprentissage a été réalisé.
\end{itemize}




\end{frame}


\section{WP Communication}


\begin{frame}
\frametitle{WP Communication} 

WP Communication: objectifs
\begin{itemize}
    \item Faciliter la coordination des autres groupes de travail.
    \item Promouvoir le projet et ses résultats auprès de potentiels utilisateurs.
\end{itemize}


\end{frame}


\begin{frame}

Les principales réalisations de ce groupe de travail sont:
\begin{itemize}
 
    \item les rapports. %CHECK ! ah ce n'est pas matthieu qui les a fait ? reponse: chefs de WP rédigent --> WP communication redaction globale --> Matthieu relecture et modification --> Enseignant
    \item cette présentation.
    \item des formations HTML et Latex. %QUOI ????
    \item le site internet.

\end{itemize}
\end{frame}



\begin{frame}
%% intérêt limité de cette liste ?
Un site internet a été réalisé et comporte:
\begin{itemize}
    \item Un forum 
    \item Un wiki
    \item Une présentation du projet 
    \item Nos résultats
    \item Une page de contacts
    \item La documentation de Coquille %NON. N'en déplaise à Eddy, il n'y a pas qu'une documentation de coquille
     % S'il s'agit de l'IDE voici encore une raison pour CHANGER DE NOM POUR L'IDE !!
\end{itemize}


\end{frame}



\section{Bibliographie}
\begin{frame}
\frametitle{Bibliographie} 

 \begin{itemize}
\item [1] http://coq.inria.fr/stdlib/ [la coq standard library]
\item [2] http://c-corn.cs.ru.nl/documentation/toc.html% [les modules disponibles dans Corn]  AH? Parce qu'on a utilisé C-Corn ?
\item [3] http://graal.ens-lyon.fr/coquille/
\item [4] http://perso.ens-lyon.fr/jeanmarie.madiot/coquille/forum % Utile ?
\item [5] http://coquille.wikispot.org/ %% OBSOLÈTE ! Et très utile ...
\item [6] Artifical intelligence, a modern approach, S. Russel P. Norvig

\end{itemize}
\end{frame}


%\section{Conclusion}
%\begin{frame}
 
%\end{frame}

\end{document}
